\documentclass[12pt]{article}

\usepackage[a4paper, margin=2.5cm]{geometry}
\usepackage{setspace}
\usepackage{hyperref}

\title{\textbf{Project Report}\\
\vspace{0.3cm}
02180 - Introduction to Artificial Intelligence\\
Technical University of Denmark}

\date{March 2026}

\begin{document}

\maketitle
\section*{Group number: 24}
\section*{Group Composition}
\begin{itemize}
    \item \textbf{Marco Cola} (Student ID: s260201)
    \item \textbf{Michele Berlato} (Student ID: s250437)
    \item \textbf{Alexandros Kouros} (Student ID: s260329)
    \item \textbf{Mattia Russo} (Student ID: s257370)
\end{itemize}

\vspace{0.3cm}
\newpage

\section{Game rules}
\textit{Mancala is a classic "sowing" game played on a board consisting of two rows of
pits and two larger reservoirs called Kalahs. The objective is straightforward: to accumulate 
more stones in your own Kalah than your opponent. On each turn, a player selects all stones 
from one of their pits and redistributes them one-by-one in a counter-clockwise direction, 
including their own Kalah but skipping the opponent's.
The game’s tactical depth arises from two key mechanics: Extra Turns, granted when the last stone 
lands in the player’s Kalah, and Captures, triggered when the last stone lands in an empty pit 
on the player’s side, allowing them to seize the stones in the opponent's opposite pit. The game 
concludes when one side of the board is cleared, at which point the player with the highest total count in their Kalah is declared the winner.}

\section{Game Type}
\textit{What kind of game is it? Relevant questions are e.g.: Single- or multi-player? Competitive
or cooperative? Zero-sum? Is it turn-based, concurrent or real-time? Perfect or imperfect
information (full or partial observability)? Deterministic or stochastic? And what does
this imply for the kind of algorithms that apply to the game?}

\section{Size space}
\textit{What, approximately, is the size of the state space of the game, that is, the number of
reachable states from an initial state of the game? Finding the precise number can often
be very difficult, but here it is enough to estimate whether it is e.g. in the ballpark of $10^5$,
$10^{10}$, $10^{15}$, $10^{20}$... You can e.g. try to give an upper bound on the size
of the state space (or both a lower and upper bound). What does this imply for the kind
of algorithms that apply to the game?}

\section{Elements of the game}
\textit{Describe the following elements of your game (or the relevant subset): $s_0$ (initial state),
Player(s), Action(s), Result(s, a), Terminal-Test(s), Utility(s, p). You don't
necessarily have to specify these formally, but try to describe them as clearly and precisely
as possible. If your chosen game cannot be expressed fully in terms of these functions,
explain why, and describe the additional components as precisely as possible.}

\section{AI states and moves representation}
\textit{How are states and moves represented in the computer? To make AI for this game, was
it enough to represent a game state as the position of the pieces or the distribution of the
cards? Or was it e.g. necessary to represent game states as belief states, and why? Don’t
give too implementation-specific answers, but answer in overall terms of algorithms and
data structures.}

\section{Methods \& Algorithms}
\textit{Which of the methods and algorithms considered in the course could potentially be used
to construct AI for your game, and which have you chosen? If you use one of the standard
algorithms without modification, you don’t have to describe it, but can just refer to the
textbook (or other sources). However, try to give a small concrete example of how the
algorithm works in your game, e.g. by illustrating a few steps of the search graph/tree
being built (you can do this in a smaller or simplified version of the game, if you prefer).
Also, argue of the appropriateness of your chosen algorithm for the chosen game.}

\section{Heuristics and Evaluation functions}
\textit{If you use heuristics or evaluation functions, make sure to explain how these have been
designed. Use mathematical notation rather than implementation-specific notation.}

\section{Implemented AI's parameters}
\textit{Are there any parameters that can be adjusted in your implemented AI? Did you try out
several different algorithms, search depths, heuristics or evaluation functions? If so, it is
a very good idea to include benchmarks (test results) comparing the different versions. It
is also interesting to compare the strength of the AIs to human players (yourselves).}

\section{Future work}
\textit{Comments on future work: How could your solution be improved? Could it be optimised
using the same algorithms and data structures that you are currently using? Or would
significant improvements of the AI require different methods, and if so, which could be
considered?}

\end{document}
